\section*{Chapter \ref{ch:b3:electrode-kinetics} --- Electrode Kinetics and Electroanalysis}

\begin{solution}{exo:b3:electrode-kinetics:1}
Cathodic slope $2.303RT/\alpha F = \qty{118}{mV}$ gives $\alpha = 59.2/118 = 0.50$; $j_0 = \qty{2.0e-6}{A.cm^{-2}}$,
the intercept at $\eta = 0$.
\end{solution}

\begin{solution}{exo:b3:electrode-kinetics:2}
$i_0 = \qty{2.0e-4}{A}$; $R_{\mathrm{ct}} = RT/(Fi_0) = 0.02569/\num{2.0e-4} = \qty{128}{\ohm}$.
\end{solution}

\begin{solution}{exo:b3:electrode-kinetics:3}
For \ce{MgSO4}, $I = \frac12(4c + 4c) = 4c$. At \qty{1.0}{mmol/L}, $I = \qty{4.0}{mol.m^{-3}}$ and $\kappa^{-1} =
\qty{0.96}{nm} \times \sqrt{100/4} = \qty{4.8}{nm}$; at \qty{0.10}{mol/L}, $I = \qty{400}{mol.m^{-3}}$, $\kappa^{-1} = \qty{0.48}{nm}$.
\end{solution}

\begin{solution}{exo:b3:electrode-kinetics:4}
$i_{\mathrm c} = C_{\mathrm{dl}}Av = \num{20e-6} \times 0.0707 \times 0.1 = \qty{0.14}{\micro A}$, under 1\,\% of the peak. At
\qty{10}{V/s} it is \qty{14}{\micro A}, while the peak grows only to $19\sqrt{100} = \qty{190}{\micro A}$:
the charging current grows as $v$, the faradaic one as $\sqrt v$.
\end{solution}

\begin{solution}{exo:b3:electrode-kinetics:5}
$D = \pi\bigl(i\sqrt t/nFAc^*\bigr)^2 = \pi(\num{12.2e-6}/(96\,485 \times 0.0707 \times \num{1.00e-6}))^2 =
\qty{1.0e-5}{cm^2.s^{-1}}$.
\end{solution}

\begin{solution}{exo:b3:electrode-kinetics:6}
$FAc^* = \qty{1.364e-2}{C.cm^{-1}}$, $F/RT = \qty{38.92}{V^{-1}}$. At \qty{50}{mV/s}: $\sqrt{38.92 \times 0.05 \times \num{7.0e-6}} =
\num{3.69e-3}$, $i_{\mathrm p} = 0.4463 \times \num{1.364e-2} \times \num{3.69e-3} = \qty{22.5}{\micro A}$; at \qty{500}{mV/s}, $\sqrt{10}$ times
more, \qty{71}{\micro A}.
\end{solution}

\begin{solution}{exo:b3:electrode-kinetics:7}
(a) Reversible. (b) \omterm{def:b3:electrode-kinetics:reversibility}{Quasi-reversible}: the separation grows with the scan rate.
(c) A chemical step follows the electron transfer (EC): at slow scans the
product is consumed before the return sweep; a fast scan outruns the chemistry.
\end{solution}

\begin{solution}{exo:b3:electrode-kinetics:8}
$K_{ij}a_j/a_i = \num{2e-4} \times 0.14/\num{4.0e-3} = 0.007$: the electrode reads 0.7\,\% high.
\end{solution}

\begin{solution}{exo:b3:electrode-kinetics:9}
Least squares: $b = \qty{0.1093}{\micro A.L.\micro g^{-1}}$, $a = \qty{0.466}{\micro A}$, $s = \qty{0.011}{\micro A}$; $c_0 = a/b =
\qty{4.26}{\micro g.L^{-1}}$, with $u(c_0) = \frac sb\sqrt{\frac1n + \frac{\bar y^2}{b^2S_{xx}}} = \qty{0.12}{\micro g.L^{-1}}$
($n = 4$, $\bar y = 1.286$, $S_{xx} = 125$).
\end{solution}

\begin{solution}{exo:b3:electrode-kinetics:10}
$Q = 0.0500 \times 1800 = \qty{90.0}{C}$; $n(\ce{Cu}) = Q/2F = \qty{4.66e-4}{mol}$, \qty{29.6}{mg}. The
result depends only on the charge, measured from a current and a time, and on
Faraday's constant: no standard is needed, provided the electrolysis is
complete and its current efficiency is 100\,\%.
\end{solution}

\begin{solution}{exo:b3:electrode-kinetics:11}
$\sqrt{DFv/RT} = \sqrt{\num{1.0e-5} \times 38.92 \times 0.1} = \qty{6.24e-3}{cm.s^{-1}}$: $\Lambda = 1.6$, \omterm{def:b3:electrode-kinetics:reversibility}{quasi-reversible}.
At \qty{10}{V/s}, $\Lambda = 0.16$, still \omterm{def:b3:electrode-kinetics:reversibility}{quasi-reversible} but further from the reversible
limit; the wave would be reversible only below about \qty{1}{mV/s}.
\end{solution}

\begin{solution}{exo:b3:electrode-kinetics:12}
$s/b = \qty{0.0838}{mM}$ and $b^2S_{xx} = \qty{203.5}{\micro A^2}$. At $\bar y$: $0.0838\sqrt{1 + 1/7} = \qty{0.090}{mM}$.
At \qty{18.0}{\micro A}: $0.0838\sqrt{1.143 + 83.0/203.5} = \qty{0.104}{mM}$. Three readings at $\bar y$:
$0.0838\sqrt{1/3 + 1/7} = \qty{0.058}{mM}$.
\end{solution}

\begin{solution}{pb:b3:electrode-kinetics:1}
\textbf{1.} $\ce{C6H12O6 -> C6H10O6 + 2 H+ + 2 e-}$ and $\ce{[Fe(CN)6]^3- + e- -> [Fe(CN)6]^4-}$;
overall
\[
  \ce{C6H12O6 + 2 [Fe(CN)6]^3- -> C6H10O6 + 2 [Fe(CN)6]^4- + 2 H+} .
\]
At the electrode, $\ce{[Fe(CN)6]^4- -> [Fe(CN)6]^3- + e-}$.
\textbf{2.} Two.
\textbf{3.} Dissolved oxygen in blood is low and variable, and its product,
hydrogen peroxide, is oxidised only at a high potential where many other
species react; the mediator is present in a known excess.
\textbf{4.} So that its surface concentration is zero: the current is then limited
by diffusion only and insensitive to small changes of potential.
\textbf{5.} The current falls as $t^{-1/2}$; readings must be taken at the time used for the
calibration.
\textbf{6.} Cottrell: $i\sqrt t = 42.0$, 41.4, 41.7, 41.6, 41.8: constant within 1\,\%.
\textbf{7.} $D = \pi(\text{slope}/FAc)^2 = \pi(\num{41.8e-6}/(96\,485 \times 0.0300 \times \num{1.00e-5}))^2 = \qty{6.55e-6}{cm^2.s^{-1}}$.
\textbf{8.} $\sqrt{\pi \times \num{6.55e-6} \times 5} = \qty{0.010}{cm} = \qty{100}{\micro m}$: the depletion reaches the top of
the layer at about \qty{5}{s}; later readings would fall below the Cottrell law.
\textbf{9.} Half: \qty{9.35}{\micro A}.
\textbf{10.} The background current: oxidation of interferents and impurities,
charging, the mediator's own residual hexacyanoferrate(II).
\textbf{11.} Yes: the residuals, at most \qty{0.095}{\micro A}, scatter without trend. At high glucose the mediator, present in
a limited amount, or the \omterm{def:b3:complex-kinetics:enzyme}{enzyme} saturates and the line bends.
\textbf{12.} $(7.10 - 0.356)/0.9189 = \qty{7.34}{mM}$.
\textbf{13.} $0.0838\sqrt{1 + 1/7 + (7.10 - 8.89)^2/203.5} = 0.0838 \times 1.076 = \qty{0.090}{mM}$.
\textbf{14.} The $1/m$ term (a single reading); repeated readings, or several strips,
reduce it.
\textbf{15.} $3 \times 0.077/0.919 = \qty{0.25}{mM}$.
\textbf{16.} $7.34 \times 180.16 = \qty{1322}{mg.L^{-1}} = \qty{132}{mg.dL^{-1}}$.
\textbf{17.} It adds its own oxidation current: the meter reads too high.
\textbf{18.} A \omterm{def:b3:electrode-kinetics:cv}{voltammogram} of the strip without glucose shows an anodic wave of
ascorbate where hexacyanoferrate(II) is oxidised.
\textbf{19.} Fewer substances are oxidised at a lower potential.
\textbf{20.} The blank current is subtracted from the \omterm{def:b3:complex-kinetics:enzyme}{enzyme} electrode's reading
before the calibration line is applied.
\textbf{21.} $D$ grows by a few per cent per kelvin and the current as $\sqrt D$: meters measure
the temperature and correct for it.
\textbf{22.} Viscosity and the red-cell fraction of blood change $D$ and the current;
standards in a similar matrix cancel these effects.
\textbf{23.} Strip-to-strip variations (electrode area, \omterm{def:b3:complex-kinetics:enzyme}{enzyme} and mediator loading),
temperature and blood composition add to the calibration's own uncertainty.
\textbf{24.} \textbf{$c(\text{glucose}) = 7.34 \pm \qty{0.09}{mmol.L^{-1}}$} (standard uncertainty), about \qty{132}{mg.dL^{-1}}.
\end{solution}
